\section[ZFS ]{\texorpdfstring{ZFS \footnote{\label{fn:ref_29}https://openzfs.org/wiki/System\_Administration}}{ZFS }}\label{zfs-1}

In ZFS si segue questo principio base: \emph{Un blocco in uso non viene
mai sovrascritto.}

Modificare dati vuol dire dunque scrivere il contenuto in una nuova
porzione di memoria, ed il filesystem dovrà essere modificato per
puntare a questa nuova area.

Questo principio, permette di implementare e gestire facilmente feature
avanzate come snapshots ed il self-healing dei dati.

In ZFS lo storage è organizzato all'interno di una pool, ossia una
struttura ad albero:

\begin{itemize}
\tightlist
\item
  le foglie contengono i dati
\item
  i nodi interiori sono dei \textbf{block pointers}
\item
  l'\textbf{uberblock} è il nodo radice, tutta la pool è navigabile a
  partire da qui.
\end{itemize}

Un block pointer tiene traccia di un insieme di blocchi e ne traccia la
grandezza e la \textbf{checksum}. Ogni nodo dell'albero calcola la
checksum in base al contenuto delle informazioni dei figli. Quindi,
modificare dati nel filesystem vuol dire ricalcolare le checksum e
ricostruire la catena dei padri a partire dalla foglia, aggiornando le
informazioni.

Dopo aver applicato la catena di modifiche a partire dalla foglia, si
ottiene un nuovo \(\text{uberblock'}\), ossia il nuovo albero che
descrive il filesystem con i nuovi dati e le checksum aggiornate.
L'intera operazione descritta fino ad ora è \textbf{atomica}.

Il costo per una modifica nella pool è proporzionale alla profondità
dell'albero, che non è necessariamente legato alla grandezza in byte
della pool. Questo approccio permette di tenere in memoria sia
\(\text{uberblock'}\) che la vecchia radice \(\text{uberblock}\).

ZFS segue un modello transazionale: le operazioni sono organizzate in
gruppi di transazioni chiamati \passthrough{\lstinline!txg!} e l'intero
gruppo viene eseguito atomicamente.

\subsection{Vantaggi di Pool e
transazioni}\label{vantaggi-di-pool-e-transazioni}

\emph{Gli snapshot del filesystem vengono eseguiti a costo quasi zero}:
fare uno snapshot vuol dire, chiedere al filesystem di non liberare i
dati delle foglie se questi devono essere eliminati o modificati in
versioni future dell'albero. Per indirizzare i dati bisogna tenere in
memoria l'\(\text{uberblock}\) e l'intero albero che descrive l'albero
al momento dello snapshot.

\begin{quote}
Uno snapshot in ZFS è una copia read-only del sistema o di un volume ad
un determinato istante. Creare uno snapshot è un'operazione quasi
istantanea, non consuma spazio ed è atomica.
\end{quote}

\textbf{ZFS è in grado di identificare gli errori}: nei blocchi grazie
al controllo della corrispettiva checksum che risiede nel padre. Se
nella pool esiste una copia, ridondante dei dati, allora ZFS può
automaticamente ripristinarne il contenuto.

\subsection{Quando ZFS diventa
costoso?}\label{quando-zfs-diventa-costoso}

Modificare un file vuol dire allocare qua e là nella pool nuovi blocchi,
rompendo la sequenzialità dei blocchi nel file system. Ne consegue che
le performance di ZFS in fase di allocazione dipendono da quanto è
grande la pool e da quanto è frammentata la memoria libera.

\begin{quote}
Se la pool è vuota allora trovare spazio contiguo è facile. Altrimenti
l'allocatore deve lavorare per trovare spazio per nuovi dati.
\end{quote}

Inoltre bisogna tenere presente che per fare qualsiasi modifica, anche
un'eliminazione, ZFS ha bisogno di memoria per creare i nuovi metadati.
In generale per fare una modifica parziale ad un record, bisogna
leggerlo, modificarlo e scriverlo dentro ad un nuovo record da attaccare
all'albero. Dunque piccole modifiche frequenti possono essere costose.
